\documentclass[12pt]{article}
\usepackage[margin=1in]{geometry}
\usepackage{enumitem}
\usepackage{titlesec}
\usepackage{parskip}
\usepackage{graphicx}
\usepackage{xcolor}
\usepackage{hyperref}
\usepackage{fontenc}

\title{\textbf{Heller and Mershon 2005 RG Notes}\\
\large Heller, William B., and Carol Mershon. ``Party Switching in the Italian Chamber of Deputies, 1996--2001'' [2005].\\
\textit{The Journal of Politics} 67(2): 536-559.}
\author{}
\date{}

\begin{document}

\maketitle

Strongly relates to the EVL reasoning from Clark et al 2017 --- party members disagreeing with their party's policy direction may exit, voice their opposition and subsequently seek to change the policy direction, or simply remain loyal to the party. This essentially outlines the conditions leading to exit as opposed to the V/L alternatives

Two primary dimensions in the decision to switch parties / ``exit'' --- electoral dimensions suggest that MPs who rely on ``personal vote coalitions'' will be less likely to switch parties \textit{if party labels are clear such that voters can easily distinguish between them. }Subsequently less salient for legislators elected primarily by means of their party membership

Italian MPs from less-ideologically-defined parties (largely centrist parties) are much more likely to party-switch than their peers in strongly (and clearly) ideological parties

Though unclear ideology may significantly associate with party switching, it does not seem to be associated with bloc switching --- that is, even when a party is not strongly ideological, its members may be expected to depart for other parties, but less likely to depart for parties outside of their own bloc. (Some of this is tautological --- all bloc switches are party switches, but not all party switches are bloc switches.)

Party switches occur significantly more frequently in the immediate lead up to a non-parliamentary election than in the months following; the reverse is true for changes in government, where switches are more likely at or shortly after the turnover.

Primary push mechanism is subsequently uncertainty / ambiguity over party identity and idelogy

\end{document}
